% ============================================================================
% CHAPTER 7: FIRST AND SECOND ORDER SYSTEM RESPONSES
% ============================================================================

\chapter{First and Second Order System Responses}
\label{ch:standard_responses}

The vast majority of physical systems encountered in control engineering can be approximated by first or second-order models. Understanding these canonical forms provides the foundation for analyzing and designing control systems across all engineering domains.

% ============================================================================
\section{First-Order Systems}
\label{sec:first_order}
% ============================================================================

\subsection{Standard Form and Parameters}

The standard first-order transfer function is:
\begin{equation}
G(s) = \frac{K}{\tau s + 1}
\label{eq:first_order_tf}
\end{equation}

where:
\begin{itemize}
    \item $K$ = DC gain (steady-state gain), dimensionless or with appropriate units
    \item $\tau$ = time constant (seconds)
\end{itemize}

The single pole is located at $s = -1/\tau$, which lies on the negative real axis for stable systems ($\tau > 0$).

\subsection{Step Response Derivation}

For a unit step input $R(s) = 1/s$, the output in the Laplace domain is:
\begin{equation}
Y(s) = G(s) \cdot R(s) = \frac{K}{\tau s + 1} \cdot \frac{1}{s} = \frac{K}{s(\tau s + 1)}
\end{equation}

Using partial fraction expansion:
\begin{equation}
Y(s) = \frac{K}{s(\tau s + 1)} = \frac{A}{s} + \frac{B}{\tau s + 1}
\end{equation}

Solving for coefficients: $A = K$ (setting $s = 0$) and $B = -K\tau$ (setting $s = -1/\tau$).

Taking the inverse Laplace transform:
\begin{equation}
\boxed{y(t) = K\left(1 - \ee^{-t/\tau}\right), \quad t \geq 0}
\label{eq:first_order_step}
\end{equation}

\begin{keyconceptbox}[title=First-Order Step Response Characteristics]
\begin{itemize}
    \item At $t = \tau$: output reaches $63.2\%$ of final value ($1 - \ee^{-1} \approx 0.632$)
    \item At $t = 2\tau$: reaches $86.5\%$ of final value
    \item At $t = 3\tau$: reaches $95.0\%$ of final value
    \item At $t = 4\tau$: reaches $98.2\%$ of final value (practical settling time)
    \item At $t = 5\tau$: reaches $99.3\%$ of final value
    \item Initial slope: $\dd y/\dd t|_{t=0} = K/\tau$
    \item Final value: $y(\infty) = K$
\end{itemize}
\end{keyconceptbox}

% ============================================================================
\section{Second-Order Systems}
\label{sec:second_order}
% ============================================================================

\subsection{Standard Form and Parameters}

The standard second-order transfer function is:
\begin{equation}
G(s) = \frac{K\omega_n^2}{s^2 + 2\zeta\omega_n s + \omega_n^2}
\label{eq:second_order_tf}
\end{equation}

where:
\begin{itemize}
    \item $K$ = DC gain (steady-state value for unit step input)
    \item $\omega_n$ = undamped natural frequency (rad/s)
    \item $\zeta$ = damping ratio (dimensionless)
\end{itemize}

\begin{definitionbox}[title=Definition: Damping Ratio]
The damping ratio $\zeta$ quantifies the level of damping relative to critical damping. It determines the qualitative nature of the system response:
\begin{itemize}
    \item $\zeta = 0$: Undamped (perpetual oscillation)
    \item $0 < \zeta < 1$: Underdamped (decaying oscillation)
    \item $\zeta = 1$: Critically damped (fastest non-oscillatory response)
    \item $\zeta > 1$: Overdamped (sluggish, non-oscillatory response)
\end{itemize}
\end{definitionbox}

\subsection{Pole Locations}

The characteristic equation $s^2 + 2\zeta\omega_n s + \omega_n^2 = 0$ yields poles at:
\begin{equation}
\boxed{s_{1,2} = -\zeta\omega_n \pm \omega_n\sqrt{\zeta^2 - 1}}
\label{eq:second_order_poles}
\end{equation}

For the underdamped case ($0 < \zeta < 1$), the poles are complex conjugates:
\begin{equation}
s_{1,2} = -\zeta\omega_n \pm \jj\omega_n\sqrt{1 - \zeta^2} = -\sigma \pm \jj\omega_d
\end{equation}

where $\sigma = \zeta\omega_n$ is the decay rate and $\omega_d = \omega_n\sqrt{1-\zeta^2}$ is the damped natural frequency.

% ============================================================================
\section{Underdamped Response ($0 < \zeta < 1$)}
\label{sec:underdamped}
% ============================================================================

\subsection{Complete Step Response Derivation}

For a unit step input, the Laplace-domain output is:
\begin{equation}
Y(s) = \frac{K\omega_n^2}{s(s^2 + 2\zeta\omega_n s + \omega_n^2)}
\end{equation}

Using partial fractions with complex poles $s = -\sigma \pm \jj\omega_d$:
\begin{equation}
Y(s) = \frac{K}{s} - K\frac{s + 2\zeta\omega_n}{s^2 + 2\zeta\omega_n s + \omega_n^2}
\end{equation}

Completing the square in the denominator and manipulating:
\begin{equation}
Y(s) = \frac{K}{s} - K\frac{(s + \sigma) + \sigma}{(s + \sigma)^2 + \omega_d^2}
\end{equation}

\begin{equation}
Y(s) = \frac{K}{s} - K\frac{s + \sigma}{(s + \sigma)^2 + \omega_d^2} - K\frac{\sigma}{\omega_d}\frac{\omega_d}{(s + \sigma)^2 + \omega_d^2}
\end{equation}

Taking the inverse Laplace transform:
\begin{equation}
y(t) = K\left[1 - \ee^{-\sigma t}\cos(\omega_d t) - \frac{\sigma}{\omega_d}\ee^{-\sigma t}\sin(\omega_d t)\right]
\end{equation}

Using the identity $\cos\theta + \frac{\zeta}{\sqrt{1-\zeta^2}}\sin\theta = \frac{1}{\sqrt{1-\zeta^2}}\sin(\theta + \phi)$ where $\phi = \arctan\left(\frac{\sqrt{1-\zeta^2}}{\zeta}\right) = \arccos(\zeta)$:

\begin{equation}
\boxed{y(t) = K\left[1 - \frac{\ee^{-\zeta\omega_n t}}{\sqrt{1-\zeta^2}}\sin(\omega_d t + \phi)\right], \quad \phi = \arccos(\zeta)}
\label{eq:underdamped_step}
\end{equation}

An equivalent form using cosine is:
\begin{equation}
y(t) = K\left[1 - \frac{\ee^{-\zeta\omega_n t}}{\sqrt{1-\zeta^2}}\cos(\omega_d t - \psi)\right], \quad \psi = \arctan\left(\frac{\zeta}{\sqrt{1-\zeta^2}}\right)
\end{equation}

% ============================================================================
\section{Critically Damped Response ($\zeta = 1$)}
\label{sec:critically_damped}
% ============================================================================

When $\zeta = 1$, the poles are real and repeated at $s = -\omega_n$:
\begin{equation}
G(s) = \frac{K\omega_n^2}{(s + \omega_n)^2}
\end{equation}

For a unit step input:
\begin{equation}
Y(s) = \frac{K\omega_n^2}{s(s + \omega_n)^2} = \frac{K}{s} - \frac{K}{s + \omega_n} - \frac{K\omega_n}{(s + \omega_n)^2}
\end{equation}

Taking the inverse Laplace transform:
\begin{equation}
\boxed{y(t) = K\left[1 - (1 + \omega_n t)\ee^{-\omega_n t}\right]}
\label{eq:critical_step}
\end{equation}

Critical damping represents the boundary between oscillatory and non-oscillatory behavior, providing the fastest approach to steady state without overshoot.

% ============================================================================
\section{Overdamped Response ($\zeta > 1$)}
\label{sec:overdamped}
% ============================================================================

When $\zeta > 1$, the poles are real and distinct:
\begin{equation}
s_1 = -\zeta\omega_n + \omega_n\sqrt{\zeta^2 - 1}, \quad s_2 = -\zeta\omega_n - \omega_n\sqrt{\zeta^2 - 1}
\end{equation}

Let $\alpha = \omega_n(\zeta - \sqrt{\zeta^2 - 1})$ and $\beta = \omega_n(\zeta + \sqrt{\zeta^2 - 1})$, so $s_1 = -\alpha$ and $s_2 = -\beta$.

The step response becomes:
\begin{equation}
\boxed{y(t) = K\left[1 - \frac{\beta}{\beta - \alpha}\ee^{-\alpha t} + \frac{\alpha}{\beta - \alpha}\ee^{-\beta t}\right]}
\label{eq:overdamped_step}
\end{equation}

The response is dominated by the slower pole (smaller $\alpha$), resulting in sluggish behavior.

% ============================================================================
\section{Pole-Zero Map and Damping Relationship}
\label{sec:pz_damping}
% ============================================================================

The location of poles in the complex $s$-plane directly determines system behavior. For second-order systems, key geometric relationships exist.

\subsection{Geometric Interpretation}

For underdamped systems, the poles lie on a circle of radius $\omega_n$ centered at the origin:
\begin{equation}
|s| = \sqrt{\sigma^2 + \omega_d^2} = \sqrt{(\zeta\omega_n)^2 + (\omega_n\sqrt{1-\zeta^2})^2} = \omega_n
\end{equation}

The angle $\theta$ from the negative real axis satisfies:
\begin{equation}
\cos\theta = \zeta, \quad \sin\theta = \sqrt{1 - \zeta^2}
\end{equation}

\begin{keyconceptbox}[title=Pole Location and System Behavior]
\begin{itemize}
    \item \textbf{Distance from origin} ($\omega_n$): Determines speed of response
    \item \textbf{Real part} ($-\sigma = -\zeta\omega_n$): Determines decay rate of oscillations
    \item \textbf{Imaginary part} ($\pm\omega_d$): Determines frequency of oscillation
    \item \textbf{Angle from negative real axis}: $\theta = \arccos(\zeta)$ determines overshoot
    \item Lines of constant $\zeta$ are radial lines from the origin
    \item Lines of constant $\omega_n$ are circular arcs centered at the origin
    \item Lines of constant $\sigma$ are vertical lines parallel to the imaginary axis
\end{itemize}
\end{keyconceptbox}

% ============================================================================
\section{Time-Domain Specifications}
\label{sec:time_specs}
% ============================================================================

Control system design often begins with time-domain specifications that must be translated into requirements on $\zeta$ and $\omega_n$.

\subsection{Rise Time}

Rise time $t_r$ is typically defined as the time to go from 10\% to 90\% of the final value. For underdamped second-order systems, an accurate approximation is:
\begin{equation}
\boxed{t_r \approx \frac{1.8}{\omega_n}} \quad \text{(for } 0.3 < \zeta < 0.8\text{)}
\end{equation}

A more precise formula valid for broader ranges:
\begin{equation}
t_r \approx \frac{1 + 1.1\zeta + 1.4\zeta^2}{\omega_n}
\end{equation}

\subsection{Peak Time}

Peak time $t_p$ is when the first overshoot occurs. Setting $\dd y/\dd t = 0$ for the underdamped response:
\begin{equation}
\frac{\dd y}{\dd t} = \frac{K\omega_n}{\sqrt{1-\zeta^2}}\ee^{-\zeta\omega_n t}\sin(\omega_d t) = 0
\end{equation}

This occurs when $\sin(\omega_d t) = 0$, giving $\omega_d t = n\pi$. The first peak is at:
\begin{equation}
\boxed{t_p = \frac{\pi}{\omega_d} = \frac{\pi}{\omega_n\sqrt{1-\zeta^2}}}
\label{eq:peak_time}
\end{equation}

\subsection{Percent Overshoot}

The maximum overshoot $M_p$ is:
\begin{equation}
M_p = y(t_p) - K = K\frac{\ee^{-\zeta\omega_n t_p}}{\sqrt{1-\zeta^2}}\sin(\omega_d t_p + \phi) - K
\end{equation}

Substituting $t_p = \pi/\omega_d$ and simplifying:
\begin{equation}
\boxed{\%OS = 100 \cdot \ee^{-\pi\zeta/\sqrt{1-\zeta^2}}}
\label{eq:overshoot}
\end{equation}

Inverting this relationship to find $\zeta$ from a desired overshoot:
\begin{equation}
\boxed{\zeta = \frac{-\ln(\%OS/100)}{\sqrt{\pi^2 + \ln^2(\%OS/100)}}}
\label{eq:zeta_from_os}
\end{equation}

\subsection{Settling Time}

Settling time $t_s$ is when the response remains within a specified percentage of the final value. Using the envelope $\ee^{-\zeta\omega_n t}$:

For 2\% criterion ($\ee^{-\zeta\omega_n t_s} = 0.02$):
\begin{equation}
\boxed{t_s \approx \frac{4}{\zeta\omega_n} = \frac{4}{\sigma}}
\label{eq:settling_2pct}
\end{equation}

For 5\% criterion:
\begin{equation}
t_s \approx \frac{3}{\zeta\omega_n}
\end{equation}

\begin{warningbox}[title=Common Mistakes in Time-Domain Calculations]
\begin{itemize}
    \item Using $\omega_n$ instead of $\omega_d$ in peak time calculations
    \item Forgetting that settling time depends on $\zeta\omega_n$, not $\omega_n$ alone
    \item Confusing rise time definitions (0-100\% vs 10-90\%)
    \item Applying underdamped formulas to overdamped systems
    \item Neglecting that overshoot formula only applies for $0 < \zeta < 1$
\end{itemize}
\end{warningbox}

% ============================================================================
\section{Worked Examples}
\label{sec:examples}
% ============================================================================

\begin{examplebox}[title=Example 7.1: First-Order Thermal System]
A thermal system has a transfer function $G(s) = \frac{50}{10s + 1}$ relating heater input (watts) to temperature change (degrees C). Find:
\begin{enumerate}[label=(\alph*)]
    \item The time constant and DC gain
    \item The response to a 2 W step input
    \item The time to reach 95\% of steady state
\end{enumerate}

\textbf{Solution:}

(a) Comparing with standard form $G(s) = K/(\tau s + 1)$:
\begin{align}
\tau &= 10 \text{ s (time constant)} \\
K &= 50 \text{ }^\circ\text{C/W (DC gain)}
\end{align}

(b) For a 2 W step input, the steady-state temperature rise is:
\begin{equation}
\Delta T_{ss} = K \cdot 2 = 100^\circ\text{C}
\end{equation}

The response is:
\begin{equation}
\Delta T(t) = 100(1 - \ee^{-t/10})^\circ\text{C}
\end{equation}

(c) At 95\% of steady state: $0.95 = 1 - \ee^{-t/10}$, giving:
\begin{equation}
t = -10\ln(0.05) = 30 \text{ s} = 3\tau
\end{equation}
\end{examplebox}

\begin{examplebox}[title=Example 7.2: Second-Order System Analysis]
A position control system has transfer function:
\begin{equation}
G(s) = \frac{100}{s^2 + 4s + 100}
\end{equation}

Find: (a) $\omega_n$, $\zeta$, and pole locations, (b) all time-domain specifications, (c) the step response expression.

\textbf{Solution:}

(a) Comparing with $G(s) = K\omega_n^2/(s^2 + 2\zeta\omega_n s + \omega_n^2)$:
\begin{align}
\omega_n^2 &= 100 \quad \Rightarrow \quad \omega_n = 10 \text{ rad/s} \\
2\zeta\omega_n &= 4 \quad \Rightarrow \quad \zeta = 0.2 \\
K &= 1
\end{align}

Poles: $s_{1,2} = -2 \pm \jj\sqrt{100 - 4} = -2 \pm \jj 9.80$

Damped frequency: $\omega_d = 10\sqrt{1 - 0.04} = 9.80$ rad/s

(b) Time-domain specifications:
\begin{align}
t_r &\approx \frac{1.8}{10} = 0.18 \text{ s} \\
t_p &= \frac{\pi}{9.80} = 0.32 \text{ s} \\
\%OS &= 100\ee^{-0.2\pi/\sqrt{1-0.04}} = 100\ee^{-0.641} = 52.7\% \\
t_s &\approx \frac{4}{0.2 \times 10} = 2 \text{ s}
\end{align}

(c) Step response ($\phi = \arccos(0.2) = 78.5^\circ = 1.37$ rad):
\begin{equation}
y(t) = 1 - \frac{\ee^{-2t}}{0.98}\sin(9.80t + 1.37)
\end{equation}
\end{examplebox}

\begin{examplebox}[title=Example 7.3: Design for Specified Overshoot]
Design a second-order system with no more than 10\% overshoot and settling time (2\%) of 0.5 seconds.

\textbf{Solution:}

From the overshoot requirement using Equation~\eqref{eq:zeta_from_os}:
\begin{equation}
\zeta = \frac{-\ln(0.10)}{\sqrt{\pi^2 + \ln^2(0.10)}} = \frac{2.303}{\sqrt{9.87 + 5.30}} = \frac{2.303}{3.89} = 0.591
\end{equation}

From the settling time requirement:
\begin{equation}
t_s = \frac{4}{\zeta\omega_n} = 0.5 \quad \Rightarrow \quad \zeta\omega_n = 8
\end{equation}

Therefore:
\begin{equation}
\omega_n = \frac{8}{0.591} = 13.5 \text{ rad/s}
\end{equation}

The required transfer function is:
\begin{equation}
G(s) = \frac{182.25}{s^2 + 16s + 182.25}
\end{equation}

Verification: $\omega_d = 13.5\sqrt{1-0.35} = 10.9$ rad/s, $t_p = 0.29$ s
\end{examplebox}

\begin{examplebox}[title=Example 7.4: Overdamped System Response]
A hydraulic actuator has transfer function:
\begin{equation}
G(s) = \frac{25}{s^2 + 10s + 25} = \frac{25}{(s+5)^2}
\end{equation}

(a) Identify the system type and parameters. (b) Find the step response.

\textbf{Solution:}

(a) Comparing with standard form: $\omega_n^2 = 25 \Rightarrow \omega_n = 5$ rad/s

$2\zeta\omega_n = 10 \Rightarrow \zeta = 1$ (critically damped)

The system has a repeated pole at $s = -5$.

(b) Using Equation~\eqref{eq:critical_step}:
\begin{equation}
y(t) = 1 - (1 + 5t)\ee^{-5t}
\end{equation}

Settling time (2\%): $t_s \approx 4/5 = 0.8$ s (using the approximation for critical damping)
\end{examplebox}

\begin{examplebox}[title=Example 7.5: Comparing Damping Cases]
Three systems have the same natural frequency $\omega_n = 4$ rad/s but different damping:
\begin{itemize}
    \item System A: $\zeta = 0.5$ (underdamped)
    \item System B: $\zeta = 1.0$ (critically damped)
    \item System C: $\zeta = 2.0$ (overdamped)
\end{itemize}

Compare their step responses.

\textbf{Solution:}

\textbf{System A} ($\zeta = 0.5$):
\begin{align}
\omega_d &= 4\sqrt{1 - 0.25} = 3.46 \text{ rad/s} \\
t_p &= \pi/3.46 = 0.91 \text{ s} \\
\%OS &= 100\ee^{-0.5\pi/0.866} = 16.3\% \\
t_s &= 4/(0.5 \times 4) = 2 \text{ s}
\end{align}

\textbf{System B} ($\zeta = 1.0$):
\begin{align}
\text{No overshoot}&, \text{ poles at } s = -4 \\
t_s &\approx 4/4 = 1 \text{ s}
\end{align}

\textbf{System C} ($\zeta = 2.0$):
\begin{align}
\text{Poles: } s_1 &= -4(2 - \sqrt{3}) = -1.07, \quad s_2 = -4(2 + \sqrt{3}) = -14.9 \\
t_s &\approx 4/1.07 = 3.7 \text{ s (dominated by slow pole)}
\end{align}

The critically damped system (B) has the fastest settling time without overshoot. The underdamped system (A) reaches steady state faster initially but overshoots. The overdamped system (C) is slowest overall.
\end{examplebox}

% ============================================================================
\section{Practical Considerations}
\label{sec:practical}
% ============================================================================

\begin{practicalbox}
Real systems deviate from ideal first and second-order models:
\begin{itemize}
    \item \textbf{Higher-order dynamics}: Most physical systems have additional poles; second-order models capture dominant behavior when other poles are far in the left half-plane.
    \item \textbf{Zeros}: Transfer function zeros affect overshoot and rise time; a zero near a pole can approximately cancel its effect.
    \item \textbf{Nonlinearities}: Saturation, dead zones, and hysteresis cause deviations from linear predictions.
    \item \textbf{Parameter uncertainty}: Actual values of $\zeta$ and $\omega_n$ may differ from design values; robust designs account for this variation.
    \item \textbf{Time delays}: Transport delays add phase lag not captured by rational transfer functions; they significantly affect stability.
\end{itemize}
\end{practicalbox}

\begin{keyconceptbox}[title=Chapter Summary: Key Takeaways]
\begin{itemize}
    \item \textbf{First-order systems} are characterized by a single time constant $\tau$; the step response is exponential with 63.2\% rise at $t = \tau$.
    \item \textbf{Second-order systems} are characterized by $\omega_n$ and $\zeta$; the damping ratio determines whether the response is underdamped, critically damped, or overdamped.
    \item \textbf{Pole locations} encode all dynamic behavior: distance from origin sets speed, angle from negative real axis sets damping.
    \item \textbf{Time-domain specifications} ($t_r$, $t_p$, $\%OS$, $t_s$) translate directly to requirements on $\zeta$ and $\omega_n$.
    \item \textbf{Design process}: Specify time-domain requirements $\rightarrow$ calculate required $\zeta$ and $\omega_n$ $\rightarrow$ place poles appropriately.
\end{itemize}
\end{keyconceptbox}

